\documentclass[a4paper,12pt]{article}

% --- Pacchetti Consigliati ---
\usepackage[utf8]{inputenc}
\usepackage[italian]{babel}
\usepackage[T1]{fontenc}
\usepackage{geometry}
\geometry{margin=2.5cm}
\usepackage{tabularx} % Per tabelle flessibili
\usepackage{booktabs} % Per linee delle tabelle più eleganti
\usepackage{enumitem} % Per personalizzare le liste

\title{Progetto di Basi di Dati}
\author{Federico Del Pup \\ Luigi Pa\textcommabelow{s}cu \\ Matteo Passador \\ Stefano Toneguzzo}
\date{\today}

\begin{document}

\maketitle

\section{Glossario dei termini}

\begin{tabularx}{\textwidth}{l X X}
\toprule
\textbf{Termine} & \textbf{Descrizione} & \textbf{Collegamenti} \\
\midrule
Pubblicazione & Elemento di bibliografia & Autore, Pubblicazione \\
Articolo su rivista & Tipo di pubblicazione & Autore, Pubblicazione \\
Articolo per conferenza & Tipo di pubblicazione & Autore, Pubblicazione \\
Libro & Tipo di pubblicazione & Autore, Pubblicazione, Editor \\
Tesi & Tipo di pubblicazione & Autore, Pubblicazione, Università \\
Autore & Persona che scrive una pubblicazione & Pubblicazione, Affiliazione \\
Affiliazione & Gruppo di autori & Autore \\
Editore & Entità che pubblica un libro & Libro \\
Università & Luogo di discussione della tesi & Tesi \\
\bottomrule
\end{tabularx}

\section{Strutturazione dei requisiti}

\subsection{Frasi di carattere generale}
Si supponga di aver collezionato il seguente insieme di requisiti per la progettazione di una base di dati relazionale per la gestione di una bibliografia. Una bibliografia è costituita da un insieme di pubblicazioni.

\subsection{Dettaglio dei requisiti}
\begin{itemize}[label=\textbullet]
    \item \textbf{Pubblicazioni:} Ogni pubblicazione è caratterizzata da un codice univoco, un titolo, un anno di pubblicazione e uno o più autori. È presente una lista di riferimenti bibliografici (citazioni) verso altre pubblicazioni della bibliografia. Ogni pubblicazione tiene traccia del numero di citazioni ricevute.
    
    \item \textbf{Articoli su rivista:} Si registrano il nome della rivista, il numero del volume, le pagine (iniziale e finale) e il numero totale di pagine.
    
    \item \textbf{Articoli per conferenza:} Si registrano nome e luogo della conferenza, anno di svolgimento (indipendente dall'anno di pubblicazione), pagine (iniziale/finale) e totale pagine.
    
    \item \textbf{Libri:} Identificati dal codice ISBN, includono l'editore e il numero totale di pagine.
    
    \item \textbf{Tesi:} Si memorizzano l'argomento, l'università di discussione e il numero totale di pagine.
    
    \item \textbf{Autori:} Di ogni autore si memorizzano nome, cognome, email, pagina Web e una o più affiliazioni.
    
    \item \textbf{Affiliazioni:} Descritte da nome, indirizzo completo (via, civico, città, nazione) e numero di telefono.
    
    \item \textbf{Editori:} Descritti da nome e indirizzo fisico.
    
    \item \textbf{Università:} Identificate univocamente dal nome, includono indirizzo e telefono.
\end{itemize}

\section{Tavola dei volumi}

\begin{tabularx}{\textwidth}{l c r}
\toprule
\textbf{Concetto} & \textbf{Tipo} & \textbf{Volume} \\
\midrule
Pubblicazione & E & 100.000 \\
Articolo su rivista & E & 20.000 \\
Articolo per conferenza & E & 15.000 \\
Libro & E & 10.000 \\
Tesi & E & 55.000 \\
Autore & E & 70.000 \\
Affiliazione & E & 500 \\
Editore & E & 400 \\
Università & E & 300 \\
Ente di ricerca & E & 200 \\
\midrule
Citazione & R & $100k \times 20$ \\
Scritto & R & $3,3 \times 100k$ \\
Appartenenza & R & $70k \times 2$ \\
Sede & R & 55.000 \\
Pubblica & R & 10.000 \\
\bottomrule
\end{tabularx}

\end{document}